problem:  
Let \( M^{n} \) be a compact, simply connected Riemannian manifold with **almost positive sectional curvature**—that is, the set  
\[
\{(p,\sigma) \in G_2(TM) : K_p(\sigma) > 0\}
\]
is open and dense in the Grassmannian bundle \( G_2(TM) \). Assume moreover that \( M \) admits a *gradient Ricci soliton* structure \((M,g,f)\) with potential function \(f\), i.e.  
\[
\operatorname{Ric}(g) + \nabla^2 f = \lambda g
\]
for some constant \( \lambda > 0 \).  

**Question:**  
Can such a manifold \( (M,g) \) exist in dimension \( n \ge 5 \)?  
Equivalently, is it possible for a compact, simply connected gradient shrinking Ricci soliton to have almost positive sectional curvature without being positively curved everywhere?  

Why is it a "good" problem:  
This problem probes an unexplored intersection between *Ricci flow geometry* and the theory of **almost positive curvature**. On one hand, compact shrinking Ricci solitons model singularities of Ricci flow and already enjoy partial curvature positivity results (e.g., nonnegative Ricci curvature and scalar curvature positivity). On the other hand, almost positive *sectional* curvature is a delicate intermediate curvature condition sitting between nonnegative and positive sectional curvature—often preserved under flow only in subtle ways.  

If an almost-positively curved, non-positively curved compact soliton existed, it would offer a genuinely new mechanism for forming singularities just below the strictly positively curved regime; conversely, a nonexistence result would reveal an unexpected *automatic curvature improvement phenomenon* under the soliton equation. Either outcome would bridge disparate areas:  
- **Geometric analysis:** Ricci solitons and curvature evolution.  
- **Global Riemannian geometry:** structure theory for almost-positively curved manifolds.  
- **Topology of high-dimensional manifolds:** implications for possible diffeomorphism types allowed under nearly positive sectional curvature.  

The statement is elegantly simple—“can an almost-positively curved gradient Ricci soliton exist?”—yet its resolution demands profound understanding of curvature pinching, elliptic estimates, and topological rigidity results. It opens a new lens through which the borderline between strict and almost curvature positivity could be viewed dynamically. This combination of simplicity, cross-field depth, and potential for paradigm-shifting insight makes it a truly *good* research problem.